% 原书第 37 页（PDF 第 50 页）
\BookPage{37}{50}
\BookHeading{清代台湾府与扬州府的关系}{清代台湾府与扬州府的关系\\[3pt]—以清代在台湾任官的扬州人为中心所作的初步探索}{尹章义*}
\BookSubheading{一、星野说：台湾隶属于扬州}

余文仪，《续修台湾府志》（乾隆三十年修，1765）卷一，封城志，开宗明义首谓（正文的第一句）：

\begin{BookQuote}
台湾，禹贡扬州之域；天文牛、女分野，星纪之次。1
\end{BookQuote}

春秋战国之际，中国的学者整合天文学、地理学和阴阳星卜之学，将天文的二十八宿说和地理的九州诸国连系，以星辰日月之变动而辨九州诸国之吉凶，形成了「星野」说。《史记》〈天官书〉、〈封禅书〉；《汉书》〈天文志〉、〈五行志〉为之发扬，是为全盛时期。《汉书》〈天文志〉谓：

\begin{BookQuote}
凡天文在图籍昭昭可知者，经星常宿…皆有州国官宫物类之象，其伏见蚤晚、邪正、存亡、虚实…阴阳之精本在地，而上发于天者也，政失于此则变见于彼，犹景之象形、乡之应声。2
\end{BookQuote}

星野说直到清末才稍戢3，故而，康熙二十三年（1684）台湾收归版图，次年(1685)，首任知府蒋毓英纂修《台湾府志》，「星野」成为重要的题议。该书卷

\vfill
\begin{BookNotes}
\BookNote{*}{作者任职于中国文化大学史学研究所。}
\BookNote{1.}{余文仪，《续修台湾府志》，页1。台湾银行经济研究室，《台湾文献丛刊》第121种本，1962，台北。以下简称《余志》。}
\BookNote{2.}{《汉书》〈天文志〉页1b（百衲本)。}
\BookNote{3.}{陈培桂，《淡水厅志》（同治十年，1871修）卷2〈封域志星野门〉谓：「史记天官书：『扬州，牛女分野。』福建属扬州，台湾属福建，淡水属台湾，分野当从同。或本西人说，当从『翼』『轸』；顾占星之术既渺，民生吏治无所关涉，存其绵蕞，既得天道远、人道迩。治厅县者，惟当尽人道，星占有变，则加儆焉可也。」台湾银行经济研究室，《台湾研究丛刊》第46种本(1956，台北)页35。所谓「民生吏治无所关涉」实即「姑且存之」之义也。}
\end{BookNotes}

% 原书第 38 页（PDF 第 51 页）
\BookPage{38}{51}
\noindent 一「星野」门谓：

古十二州之域所以纪地，而十二次之躔所以纪天。二十八宿…各有定位。后人因地占星，即因星辨地，遂有分野之说…扬州之域，东南至海，属于牛女，为吴越之分…八闽界在瓯粤之间，原非古扬州境，而地尽东南，遂附隶扬州…台湾远隔大海，番彝荒岛，不入职方，分野之辨，未有定指。然晋天文志，十二次始角亢，以东方苍龙为之首也。唐十二次始女虚，以十二支子为之首也。按考台湾地势，极于南而迤于东，计其道里，当在女虚之交，为南纪之极，亦当附于扬州之境，以彰一统之盛焉。4

同时的首任诸罗知县无锡进士季麒光在〈客问〉一文中也说:「（台湾）北凭瓯闽，西距交广，属扬州之分，隶女虚之躔。」高拱干，续修《台湾府志》（乾隆三十四年，1695），在「星野」门中则主张「在地为扬州，在天为牛女」。云：台自破荒以来，不载版图、不登太史；星野分属，何从而辨？然台系于闽，星野宜从闽。即以闽稽之，福建「禹贡」扬州之域，天文牛、女分野…有以台郡分野，当在女、虚之交者。虚，元枵之次，在子之辰；以台之稍迤而东，遽疑其越次、越辰，是亦坐井观天之见也。6

康熙五十六年（1717）陈梦林修《诸罗县志》则认为：

俗以附漳泉分野；不知漳泉系丑地牛女之分，与巳位无涉；故占验之家，概以台湾同岛彝一体测验，而以里差详之，台湾为翼九度。诸罗在台北角木，远照斜插，隐见属翼八度七分之三。凤山在台南，亦同属翼九度杪。牛女星纪，约略之见，不可从也。7

陈梦林虽然认为高拱干牛女之说「不可从」，可是，他对于分野之说的「言人人殊」也无可奈何：

\vfill
\begin{BookNotes}
\BookNote{4.}{蒋毓英，《台湾府志》卷1〈星野志〉，页3b-4a，中华书局影刊《台湾府志三种》，1985，北京。}
\BookNote{5.}{《余志》，<艺文志〉四，骈体，页852-853。}
\BookNote{6.}{高拱干，《台湾府志》卷1〈封城志星野门〉，页1b-3a，中华书局影刊，《台湾府志三种》本。}
\BookNote{7.}{周钟瑄、陈梦林等，《诸罗县志》页31。台湾银行经济研究室，《台湾文献丛刊》第55种本，1958，台北。}
\end{BookNotes}

% 原书第 39 页（PDF 第 52 页）
\BookPage{39}{52}
\begin{BookQuote}
考步之家，言人人殊，裨邻次舍，各有差谬；此以为轸，彼以为角；甲以为氐，乙以为房。南斗六星即斗牛之斗，而分野在北。北斗天枢，在张宿十度，而分野在南。盖自昔已然。天道幽远，其理固有不可究诘者，姑合诸说而并存之，以俟博雅知星者论定云。8
\end{BookQuote}

乾隆三十年，余文仪《续修台湾府志》又综合众家之说，详细讨论之后，得到本文篇首所引述的结论。有清一代，凡通志一；府志五（福建通志台湾府二）；厅县志十三；志稿、志略、采访册等三十又四。「牛女论」和「女虚论」聚讼纷纭，但「台湾隶属扬州说」则无异议。

星野说以天象验区域之人事，将本不相涉之事物，牵强附会，固说者渺茫，言人人殊，然以后世之行政区划、归属，于远古之九州说，扬州位居东南，故而其东与南皆属扬州之域也。

\BookSubheading{二、清初台湾府与扬州府的关系}

康熙二十三年（1684）台湾收归版图，设一府三县。蒋毓英《台湾府志》（1685完稿）云：

\begin{BookQuote}
台郡三邑之人民，计之共一万六千余丁，不及内地一小邑之户口。10
\end{BookQuote}

扬州则为三千年名城，富甲天下之通都大邑。

康熙三十六年（1697）杭州人郁永河至淡水（今台北北投）探硫，其探硫日记《裨海纪游》中，记录当时之台湾（今台南市）云：

\begin{BookQuote}
近者，海内恒苦贫，斗米百钱，民多饥色；贾人责负声，日沸阇暗；台郡独似富庶，市中百物价倍，购者无吝色，贸易之肆，期约不愆；佣人计日百钱，趋趄不应召；屠儿牧竖，腰缠常数十金，每遇摴蒱，浪弃一掷间，意不甚惜；余
\end{BookQuote}

\vfill
\begin{BookNotes}
\BookNote{8.}{同前，页320。}
\BookNote{9.}{尹章义，（台湾地方志的数量、品质与方志学的发展—《台湾地方志总目录》试析〉，《方志学理论与战后方志纂修实务国际学术研讨会》，1-95页，国史馆台湾文献馆，民国九十七年五月出版，南投市，台湾。}
\BookNote{10.}{蒋毓英，《台湾府志》卷5〈风俗志〉，页2a-3a。}
\end{BookNotes}

% 原书第 40 页（PDF 第 53 页）
\BookPage{40}{53}
\begin{BookQuote}
颇怪之。因留台久，始得其故…兵饷…植蔗为糖…米、谷、麻、豆、鹿皮、鹿脯，运之四方…台土宜稼，收获倍蓰…为贾贩通外洋诸国，则财用不匮。民富土沃，又当四达之海，即今内地民人，襁至而辐辏，皆愿出于其市。萑苻陆梁，孰不欲掩而有之？11
\end{BookQuote}

《蒋志》撰稿时，大战方戢，清廷将郑氏军民眷属安插内地，台湾空虚，人户稀微；首任诸罗知县无锡进士季麒光颇为忧心，在〈条陈台湾事宜〉文中胪列要事六项，其一为备兵饷，其二即为「招集丁民」，其文曰：

\begin{BookQuote}
国家根本之计，莫先于户口；故庶而后富为从来起化之原。台湾自伪郑归诚之后，难民丁去之，闲散丁去之，官属兵卒又去之，卑县设法招徕…南北草地，一望荒芜，得人开垦，可成沃壤。12
\end{BookQuote}

《诸罗县志》〈列传〉仅两人入传，除季麒光外，张〔王尹〕亦因「招徕垦辟，抚绥多方，流民归者如市」入传13。

经过最初入台之官吏们努力招徕以及台湾地大、土沃、需人开垦诸诱因，移民大量东渡，故而，郁永河于十二年后到达台湾探硫，先到府城（今台南）备办探硫所需物资，目睹的已是糖、米、麻、鹿等生产、贸易鼎盛，人民富足安乐，且有大量驻台兵丁之兵饷挹注的乐土。

康熙中期之招徕垦辟，尚有一重人事之障碍，即靖海侯施琅厌恶惠潮之客家人，《台海使槎录》引〈理台末议〉云：

\begin{BookQuote}
台湾始入版图，为五方杂处之区，而闽粤之人尤多。先时，郑逆窃踞海上，开垦十无二三。迨郑逆平后…终将军施琅之世，严禁粤中惠、潮之民不许渡台。盖恶惠、潮之地素为海盗渊薮，而积习未忘也。琅殁，渐弛其禁，惠、潮民乃
\end{BookQuote}

\vfill
\begin{BookNotes}
\BookNote{11.}{郁永河，《裨海纪游》，页30-31，方豪校刊，台湾银行经济研究室，《台湾文献丛刊》第44种本，1959，台北。}
\BookNote{12.}{陈文达等，《台湾县志》卷10〈艺文志〉页227-228，台湾银行经济研究室，《台湾文献丛刊》第103种本，1961，台北。}
\BookNote{13.}{周钟瑄、陈梦林等，《诸罗县志》，页49。}
\end{BookNotes}

% 原书第 41 页（PDF 第 54 页）
\BookPage{41}{54}
得越渡。14

闽粤沿海之民，善良者为渔民、为海商、为水军，若遇时变，随时可转为海盗，不仅惠潮之民然。故施琅之严禁惠潮之民东渡，必有他因。

施琅死于康熙三十五年（1696），此后，大量客家人来台，与闽人竞垦15，康熙五十年（1711），台湾知府周元文指出：

\begin{BookQuote}
闽、广之梯航日众，综稽簿籍，每岁以十数万计。16
\end{BookQuote}

遂以今之台南为中心，南及高雄、屏东，北向嘉义云林，甚至入垦台北、台中盆地17。

康熙六十年（1721）朱一贵事件爆发，南澳总兵蓝廷珍奉命东征，从弟蓝鼎元任记室，「深谙全台地理情形」18。蓝鼎元在〈与制军再论筑城书〉中谓：

\begin{BookQuote}
台湾…沃野千里，粮糈足食，舟楫之利通天下。19
\end{BookQuote}

在〈覆制军台湾经理书〉中谓：

\begin{BookQuote}
国家初设郡县，管辖不过百里，距今未四十年，而开垦流移之众，延袤二千余里，糖谷之利甲天下。20
\end{BookQuote}

台湾的生产力和消费力都很强，既庶且富，两岸商贸因而繁盛，「上通江浙、

\vfill
\begin{BookNotes}
\BookNote{14.}{黄叔璥，《台海使槎录》，页93（雍正元年，1723纂)，台湾银行经济研究室，《台湾文献丛刊》第4种本，1957，台北。}
\BookNote{15.}{尹章义，〈台湾移民开发史上与客家人相关的几个谜题〉，《中国海洋开发史论文集（第四辑）》中央研究院中山人文社会科学研究所，1911，台北。收于《台湾客家史研究》，页1-29。台北市政府客家事务委员会，2003，台北。}
\BookNote{16.}{周元文，〈申请严禁偷贩米谷详稿〉，《续修台湾府志（1712周元文修）》，页121。《中华大典台湾丛书》本，国防研究院，1968，台北。}
\BookNote{17.}{尹章义，〈台北平原拓垦史研究（1697-1792)〉，《台北文献》53/54期合刊，页1-190，台北市文献会，1981，台北。又收录于《台湾开发史研究》，联经出版公司，1989，台北。}
\BookNote{18.}{蓝鼎元，《东征集》页4，（蓝廷珍序)，台湾银行经济研究室，《台湾文献丛刊》第12种本，1958，台北。}
\BookNote{19.}{前书，页29-30。}
\BookNote{20.}{前书，页34。}
\end{BookNotes}

% 原书第 42 页（PDF 第 55 页）
\BookPage{42}{55}
\noindent 下抵闽广；来往商艘，岁至以数千计」21。朱一贵事件之后，奉命到台湾实地查勘，了解台湾状况的首任巡台御史黄叔璥，返京之后以《台海使槎录》一书复命，其书前四卷记汉人事迹，名曰《赤崁笔谈》，后四卷记原住民事迹，名曰《番俗六考》。《赤崁笔谈》卷二〈商贩〉条谓：

\begin{BookQuote}
海船多漳、泉商贾。贸易于漳州，则载丝线、漳纱、翦绒、纸料、烟、布、草席、砖瓦、小杉料、鼎铛、雨伞、柑、柚、青果、橘饼、柿饼；泉州则载磁器、纸张；兴化则载杉板、砖瓦；福州则载大小杉料、干笋、香菰；建甯则载茶。回时载米、麦、菽、豆、黑白糖、锡。番薯、鹿肉售于厦门诸海口；或载糖、靛、鱼翅至上海；小艇拨运姑苏行市。船回则载布匹、纱、缎、枲绵、凉暖帽子、牛油、金腿、包酒、惠泉酒；至浙江则载绫罗、绵、绸、ケ、纱、湖帕、绒线；甯波则载绵花、草席；至山东贩卖粗细碗碟、杉枋、糖、纸、胡椒、苏木。回日则载白蜡、紫草、药材、茧绸、麦、豆、盐、肉、红枣、核桃、柿饼；关东贩卖乌茶、黄茶、绸缎、布匹、碗、纸、糖、曲、胡椒、苏木。回日则载药材、瓜子、松子、榛子、海参、银鱼、蛏干。海壖弹丸，商旅辐辏，器物流通，实有资于内地。22
\end{BookQuote}

\begin{BookQuote}
《赤崁笔谈》卷三〈物产〉「蔗」条云：煎糖…色赤而松者，于苏州贩卖，若糖湿色黑，于上海、宁波、镇江诸处行
\end{BookQuote}

销。23

黄叔璥记录的各地特产近百种，两岸商贸市场则北起关东，南迄闽、粤，正如蓝鼎元所谓台湾「舟楫之利通天下」，在长江三角洲则有宁波、苏州、镇江、上海等地；上海、苏州与扬州相去不远；镇江于扬州更是近在密迩，扬州与台湾之商贸关系理应相当密切，郁永河、蓝鼎元、黄叔璥等人均未述及扬州，则不知何故？是郁、蓝、黄等人不愿突显扬州？抑或扬州属于消费型城市，并无特产可销台湾？或者，实际上扬州与台湾商贸往来密切而史料无征？或许扬州地方有更原始、更细致

\vfill
\begin{BookNotes}
\BookNote{21.}{陈文达等，《台湾县志》〈海道〉，页67。（康熙五十九年，1920修），台湾银行经济研究室，《台湾文献丛刊》第103种本，1961，台北。}
\BookNote{22.}{《台海使槎录》，页47-48。}
\BookNote{23.}{前书，页57。}
\end{BookNotes}

% 原书第 43 页（PDF 第 56 页）
\BookPage{43}{56}
\noindent 的史料可供解答。

\BookSubheading{三、林爽文事件两淮商人捐赀两百万}

明天启二年（1622）以前，台湾是中国沿海商人的逋逃渊薮、海盗巢穴和贸易基地；永历十五年（清顺治十八年，康熙元年，西元1661）郑成功东征，赶走荷兰人，建立隶属于大明永历帝的延平王国24；康熙二十二年（1683）施琅东征，次年设台湾一府，台湾、诸罗、凤山三县，此后，大量闽、粤人民入垦台湾，台湾逐渐迈向中土儒汉化之途，直至清末，基本上仍然未脱「移垦社会」习气，「民风剽悍、桀骜难治」。约略估计，清代治台（1683-1895）期间，「民变」（以政府为对象的反叛事件、攻打衙署和营盘）达七十余次，民间「械斗」（闽/粤，漳/泉…分籍、分类）则不计其数。

乾隆五十一年（1786）年底，台湾爆发了林爽文事件，由于天地会居中串联、运作，全台响应25。闽浙总督常青东征失利；改派陕甘总督福康安、侍卫大臣海兰察率侍卫巴图鲁一百二十大员、精锐九千东征，五十二年（1787）十月二十八自福建蚶江崇武放洋、十一月二日在鹿港登岸，十二月二十五日攻入林爽文巢穴大里杙（今台中县大里市），五十三年（1788）年正月初，擒林爽文，三月中擒庄大田于琅峤（今屏东县恒春市），乱平26。

平定林爽文事变，征集了十四省兵力，征米一百二十万石，费银千万两，天下骚动。五十二年（1787）十一月，广东洋商潘文严等捐银三十万两以充军需，盐商

\vfill
\begin{BookNotes}
\BookNote{24.}{尹章义，〈延平王国的性质及其在国史上的地位—兼答厦门大学邓孔昭教授〉，「郑成功驱荷复台三百四十周年学术研讨会」大会宣读论文，2003，厦门大学，厦门。又刊于《长共海涛论延平纪念郑成功驱荷复台340周年学术研讨会论文集》，页140-156，上海古籍出版社，2002年8月，上海。}
\BookNote{25.}{尹章义，〈林爽文事件天地会的浮现与台北第一次大械斗—台北土城大墓公的起源〉，《莲缘》创刊号，页64-76，青莲文化事业公司，2010年9月，台北。}
\BookNote{26.}{刘妮玲，《清代台湾民变研究》第五章〈民变性质之综合分析〉页266-340。台湾师范大学历史研究所，中华民国七十二年，台北。}
\end{BookNotes}

% 原书第 44 页（PDF 第 57 页）
\BookPage{44}{57}
\noindent 李念德等请捐饷银二十万两27；五十三年（1788）正月两淮盐商江广达、程俭德等「情愿公捐银二百万两，以备犒赏之需」28；同时，长芦盐商王德宜、查奕茂等捐银三十五万两，山东盐商江平庆、刘占安等捐银一十五万两29；两浙盐商何永和、方瑞基、汪有基、吴志仁等捐银四十万两，「以为赏恤兵丁之用」30，次月，又捐银三十万两，「以为添备城工之用」31。

各地商人的捐款（极可能是摊派），朝廷当然照单全收。各地商人之捐款，集中在乾隆五十二年年底、五十三年年初，显然曾经预先协商，金额的大小，似乎也经过分配；两淮盐商捐银两百万两，超过广东、两浙、长芦、山东商人之总数一百七十万两，可见两淮盐商之财力雄厚，富甲天下。

正如本文所引述郁永河、蓝鼎元和黄叔璥的记录，台湾虽然从康熙中期开始才有大量闽、粤移民入台垦荒，短短二、三十年间，已经是「糖谷之利甲天下」「舟楫之利通天下」，糖、米、麻、豆行销天下，加上国际贸易和兵饷挹注，可以说是既庶且富；虽然「民风剽悍、桀鳌难治」，械斗，民变频仍，对于台湾的总体经济发展影响似乎并不太大。因此，各地商人的捐献，对于国家财政挹注兼示忠诚的意义远大于对于台湾地区的救恤。

道光、咸丰年间（1825-1860）台湾的樟脑事业日渐兴盛；探脑之后继之植茶，同治年间（1861-1874），台茶外销大盛。光绪初，沈葆桢、王凯泰（扬州人、福建巡抚）、丁日昌、吴赞诚、岑毓英、刘铭传相继在台湾推动现代化的事业，光绪十三年（1887）台湾建省，台湾成为中国唯一出超的省份，也是最现代化的省

\vfill
\begin{BookNotes}
\BookNote{27.}{《清宫谕旨档台湾史料（一）》，（国立故宫博物院影刊本，1996，台北），乾隆五十二年十一月二十九日上谕，页822-823。(以下简称《谕旨档》)《清实录台湾史资料专辑》（福建人民出版社，1993，福州）乾隆五十二年十一月二十九条，仅录潘文严捐银三十万两，页465。(以下简称《实录》)}
\BookNote{28.}{《清宫宫中档奏折台湾史料（十）》，（国立故宫博物院影刊本，2004，台北）。乾隆五十三年正月初四全德奏，页168。(以下简称《宫中档》)《实录》乾隆五十三年一月十九日条，页498。}
\BookNote{29.}{《宫中档》页171-172，乾隆五十三年一月四日，穆腾额奏。《实录》乾隆五十三年一月六日条，页486。}
\BookNote{30.}{《宫中档》页174，乾隆五十三年一月五日，浙江巡抚觉罗琅玕奏。《实录》乾隆五十三年一月十四日条，页493。}
\BookNote{31.}{《宫中档》页252，乾隆五十三年二月二日，浙江巡抚觉罗琅玕奏。}
\end{BookNotes}

% 原书第 45 页（PDF 第 58 页）
\BookPage{45}{58}
\noindent 份32。台湾人的经济实力，可能与扬州不相上下。

光绪三、四年（1877-1878），华北苦旱，山西、河南、直隶尤甚，各地官员、商人捐款济助，台北板桥林家由林维源代表，捐银五十万元，光绪五年四月即已缴清33，不但皇上认为「晋、豫灾民赖以全活甚众」34，时人也誉之「罄私室之力，纾难公家，捐资之巨，为亘古史册所未见」35，直隶总督李鸿章、闽浙总督何璟、办理各国事务大臣丁日昌等奏称：

\begin{BookQuote}
该绅林维源捐数至五十万元之多，非捐例所能核办…此项捐数甚钜，应得叙奖，实无成案可援。36
\end{BookQuote}

所谓「非捐例所能核办」；所谓「实无成案可援」，意谓一人独捐五十万两者，前所未有也，故台湾之富，天下咸知也。

光绪五年，山西又遭亢旱而直隶却遭水灾，「苏、扬、浙、沪乐善诸君子攘臂一呼，捐集二十余万金」，「晋赈公所」之劝募函即谓「台湾富庶、海内艳称」，期待「台北绅富」能慷慨解囊37。

光绪五年，苏州、上海、浙江和扬州的商人，捐集了二十余万金；台北的板桥林维源就捐了五十万元，光绪五年四月即已缴清，可见台湾人财力的雄厚，不在苏、扬、浙、沪商人之下。

\BookSubheading{四、清代在台湾任官的扬州人初探}

中国古代职官之任用，除了教职官之外，基本上都有回避本籍的制度，亦即本地人不得在本地为官。故而，扬州人任官，都得到外省，其中一部分人派到福建，

\vfill
\begin{BookNotes}
\BookNote{32.}{尹章义，《台湾发展史》第九章〈列强再度竞逐与清廷的应对〉，页101-110。}
\BookNote{33.}{〈闽浙总督何等奏绅士特捐款全数缴清折〉，《清季申报台湾纪事辑录》页853，光绪五年四月二十三日条。台湾银行经济研究室，《台湾文献丛刊》第247种本。(以下简称《申报》)}
\BookNote{34.}{《申报》页884，光绪五年七月二十七日条〈恭录谕旨（七月十九日上谕）>}
\BookNote{35.}{〈全台协募晋赈公所致林时甫观察书〉，《申报》页920-921，光绪五年十月十四日条。}
\BookNote{36.}{〈直督李等奏绅士捐输巨款全数缴清遵旨拟请破格优奖折〉，《申报》页877-878，光绪五年七月二十一日条。}
\BookNote{37.}{台礼附登〉，《申报》页927，光绪五年十二月初六日条。}
\end{BookNotes}

% 原书第 46 页（PDF 第 59 页）
\BookPage{46}{59}
\noindent 其中又有一部分再派到台湾。本节以在台湾任官的扬州人为主，部分在福建任官，掌理台湾事务或担任闽浙总督、福建巡抚的扬州人也纳入。所谓扬州人，则以清代扬州府辖下州县籍民称之。

清初，扬州府属十州县，凡高邮、泰州、通州三州；江都、仪征、宝应、兴化、如皋、泰兴、海门七县。康熙十一年（1672）年，海门县治圮于海，裁县入通州；雍正二年（1724）通州升直隶州，如皋、泰兴二县归之。雍正九年（1731）分江都县北为甘泉县，乾隆三十三年（1768）自泰州分设东台县，民国元年（1912）又废甘泉归并于江都县。故而本文搜罗过程，即以前述十二州县为对象，并不因当今扬州市辖下或通州及直隶、甘泉归治而剔除38。实际所得二十五人，分属通州（五人）、高邮（二人）、宝应（二人）、江都（三人）、仪征（四人）、甘泉（六人）、如皋（二人）、泰州（一人）等。兹分述如下（前文已有详注之方志资料，以下探随文注）：

\BookSubheading{（一）王兆升}

高拱干《台湾府志》卷三「秩官」「台湾县知县」（页5b）:「王兆升:江南通州人，丁酉举人，康熙二十七年任，三十年升兵部职方主事。」卷十「艺文」录其「郊行即事诗」八首：

其一：

\begin{BookQuote}
奉命筹军国，非关玩物华；新凉犹未至，余暑正方赊。鸣骑依残渡，行旌带晚霞；无劳呵殿急，恐警野人家。
\end{BookQuote}

\vfill
\begin{BookNotes}
\BookNote{38.}{本文检阅以下扬州方志：1.雷应元等修，《扬州府志》，康熙三年（1664)刻本，《稀见中国地方志汇刊》影刊本。2.尹会一、程梦星等修，《扬州府志》，雍正十一年（1733）刻本，成文出版公司《中国方志丛书》影刊本，1974，台北。}
\BookNote{3.}{阿克当阿、姚文田等，《重修扬州府志》，嘉庆十五年（1810）刻本，2006年、广陵书社影刊本、扬州。}
\BookNote{4.}{英杰、晏瑞书等，《续纂扬州府志》，同治十三年刻本（1874），成文出版公司《中国方志丛书》影刊本。文中不另行详注版本。}
\end{BookNotes}

% 原书第 47 页（PDF 第 60 页）
\BookPage{47}{60}
\noindent 其二：夙笃因忧饷，三农兼课工；高资雨茂，禾好卜年丰。社鼓挝须习，豳笙调预融。麦蚕增喜笑，亦复学张公。其三：危桥恣蚁渡，蓬室类蜗房。草诧胭脂紫，花闻月下香(原注：胭脂花与月下香，皆海外异卉）。番榴生碍路，野鸠语谁乡！触目殊风景，车轮几断肠。其四：但歌灵雨降，瘠里仍然贫。早谷先垂颖，新秧犹压苹；谁怜蓣作饭（原注：俗呼土瓜；无粮代饭者十之五)。，果见釜生尘！不历郊原过，焉知民苦辛？其五：欲悉田家事，无将上下睽。儿童擎绿酒，野老馈黄梨；忘分辞方尽，陈情意转凄。莫云民俗薄，绥靖念残黎。其六：寄身重海外，嘿坐计生平；幸隶芙蓉署，欣游细柳营；深村忘犬警，远渚有蛙更；夜雨晨方歇，驱车再问程！其七：疮痍才息后，鹰眼未全。箸为军糈借，心因民瘼忉；催科良媿拙，抚字敢矜劳。珍重斯民意，输租谢尔曹！其八：飞蝗来凤邑，忽蔽永丰云；妇子奔号急，微民稽拜殷。弭灾全藉德，祝蜡岂徒文！赖有廉仁宪，驱山息异氛。陈文达《台湾县志》〈秩官志〉如之，〈艺文志〉仅收录第一首。刘良璧《重


% 原书第 48 页（PDF 第 61 页）
\BookPage{48}{61}
\noindent 修台湾府志》〈艺文志〉则收录八首39。

\BookSubheading{（二）卓尔坛}

\noindent 刘良璧《重修福建台湾府志》卷十二〈职官·武职·台湾水师协镇左营游击》卓尔坛条（页395）：卓尔坛：江都人。（章义按：年份在康熙三十二年之后两任)

\BookSubheading{（三）张骇}

\noindent 前志「台湾水师协标左营守备」张骇条（页407）：张骇：江南江都人，行伍。康熙四十七年任。又，「澎湖水师协右营游击」张骇条（页399）:张駄：江南江都人，行伍。康熙五十八年任。

\BookSubheading{(四）王振}

\noindent 前志「台湾水师协镇左营游击」王振条：（页396）：王振：江南通州人，行伍。雍正五年任。

\BookSubheading{（五）夏之芳}

\noindent 刘良璧《重修福建台湾府志》卷十三〈职官志（一）文职〉「钦命巡视台湾御史」夏之芳条（页351）：夏之芳：江南高邮人，癸卯进士。雍正六年任；留任一年。岁、科两试，校士公明，着有《海天玉尺编》二集。又，卷十五〈名宦志〉夏之芳条（页432）：夏之芳，字荔园，号筠庄；江南高邮州人。才识渊超，英明简重。登癸卯恩科进士，充南书房，以御试第一入史馆。寻转谏垣。雍正六年，巡视

\vfill
\begin{BookNotes}
\BookNote{39.}{刘良璧等，《重修福建台湾府志》（乾隆七年，1742)，台湾银行经济研究室《台湾文献丛刊》七十四种本，1961，台北。(以下凡新引书注明版本，若为前文已注明者，仅著书名。)}
\end{BookNotes}

% 原书第 49 页（PDF 第 62 页）
\BookPage{49}{62}
\noindent 台湾兼学政。至则以澄叙官方、振兴文教为己任。主岁、科两试，公慎明敏；栽培士类，悉本真诚。按巡南北二路，鸡犬不惊，民番咸悦。生平廉介，苞苴永杜；而接物复乐易冲和，绝少崖岸，尤为寮属所敬服。着有「海天玉尺」二编，至今学者宗之。

又，卷二十〈艺文志〉收录二编之两篇序文：

甲、〈海天玉尺编初集序〉:（页532-533)台湾僻处海隅，自入版图历今垂五十载。旧制：郡邑守令外，以观察使领之。康熙六十年，始以台臣出司巡视，满、汉各一员，岁一更替；盖重海疆也。我皇上御极以来，文德覃敷、声教暨讫，休风雅化，已遍遐陬；犹于临轩遣策，加意遴选，再三训诫，俾奉命诸臣得所遵循，以为观风训俗之轨，则圣明洵周且至矣。至台郡学政，旧属之观察使，不关巡视之责。丁未冬，始奉命令汉御史兼摄；睿旨初下，余适恭膺简命，出巡兹土，缘奏请学篆兼使节以行。于戊申之二月，得抵台署焉。夫台湾，山海秀结之区也。万派汪洋、一岛孤峙，磅礴郁积之气互绝千里，灵异所萃，人士必有钟其秀者。况数年来，沐国家休养教育之泽，涵濡日深，久道化成；固已家弦户诵，蒸蒸然共跻于声名文物矣。第四民之众，士为之倡；士习之邪正，风俗因之。台郡人文蔚起，宁患无才？有才不醇，则庞杂与卑污同病。昔人谓：「士先器识，而后文艺」。士习不端，祗以文藻夸世，匪唯无益，抑且为民害焉。稽台郡初辟时，岁、科抡才，多借资于漳、泉内郡。近已诏下厘剔，非生长台地者不得隶于台学。此又盛朝作养边陲之至意。都人士既得〔灵〕秀于山海钟毓，尤当厚自鼓舞，以上副圣恩，毋自域于棫朴菁莪外也。柳州云：『报国恩，唯文章』；士生此昌明之世，诚读书绩学、修身立品，使文章积为有用；而又以其诗书弦诵驯其子弟、化导乡人，俾淳庞和气，遍于蛮天菁岭间。则上以鼓吹休明、下以转移风俗，是固宣铎者所厚望；而观风训俗之责，亦可借此以仰报天子矣。兹因岁试告竣，择其文尤雅驯者付之梓，而因以发之，益使台之人知录其文者之非徒以文示也。是为序。


% 原书第 50 页（PDF 第 63 页）
\BookPage{50}{63}
\noindent 乙、〈海天玉尺编二集序〉:（页534-535）岁试既竣，择其文之拔前茅者录付剞劂，亦为海隅人士作其气而导之先路也。嗣以出巡南北，鹿鹿驱驰，且复瓜期将届；科试，拟将待之来哲矣。迨已酉正月，复奉恩纶留任，乃得于春三月举行科试事。时台郡生童忠雀跃，应试者几倍岁试之半。爰秉公莅事，慎终如始；不一月而告竣焉。台地越在海表，才隽之士，时时间出；所虑无老师宿学，穷经嗜古而陶冶之。其抱守椠铅者，甚以僻陋寡闻，销磨其志于蚓窍蛙鸣之内。才以地限，殊可惜也。然余屡试校阅，皆随材甄别，曲示鼓励，故其文亦颇渐次有可观者。大约文人之心，类从其地之风气。台士之文多旷放，各写胸臆，不能悉就准绳。其间云垂海立、鳌掣鲸吞者，应得山水奇气；又或幽岩峭壁、翠竹苍藤，雅有尘外高致。其一瓣、一香、一波、一皱，清音古响以发自然，则又得曲岛孤屿之零烟滴翠也。海天景气绝殊，故发之于文颇能各逞瑰异。至垂绅搢笏、庙堂黼黻之器，则往往鲜焉。固其士之少所涵育，亦其地之风气僻远而然也。故岁试所录，强半灵秀之篇；科试则多取醇正昌博者，为台人更进一格。亦俾知盛朝文教之隆、设科取士之法，以明正大为宗，而不得囿于方隅闻见间也乃更合岁、科试文得八十首付之梓，以为多士式。又录其〈台湾纪巡诗〉百首中之三十九首（页589-592）:海天咨度报皇华，蚤卜升平遍水涯。载德行旌才出郭，暖烟晴竹巳家家。节旄高插引晴岚，人拥花骢弛辔衔；拜罢耆童回道左，纷来朱履又青衫。野田晴晓碧天空，地指扶桑东复东；赤嵌城边云散彩，拓开海日一轮红。负暄童叟爱冬温，红稻成堆拥华门；桐竹周遭鸡犬静，教人历历认花村。牢拴竹篾怕身肥，带孔频频减旧围；爱把细腰谐凤卜，楚王宫里梦双飞。牛车无日不当官，没字邮符颠倒看；踏水冲泥何限苦，忍教横挞更无端。锄田捕鹿洽婚姻，乐事相寻满社春；嚼得瓮头姑待酒，木瓢椰碗竞奇麟。手制云箫别有腔，吹来鼻息爱成双；月明引得风前凤，未许当门夜吠厖。不须挑逗费闲心，竹片沿丝巧作琴；远韵低微传齿颊，依稀私语夜来深。


% 原书第 51 页（PDF 第 64 页）
\BookPage{51}{64}
\noindent 男拔髭须女绣颐，乍逢鉴貌尽多疑；雕题凿齿徒矜尚，未解双蛾夜画眉。杵臼轻敲似远砧，小鬟三五夜深深；可怜时办晨炊米，云磬霜钟咽竹林。秋尽官催认饷忙，一丝一粟尽输将；最怜番俗须重译，溪壑终疑饱社商。虚滩水落涨沙泥，南北中分虎尾溪；一带草荒村舍少，年来新集有蒸藜。南山中断北山连，逗漏云间半线天；道是孤城还少郭，竹环廛市起炊烟。献芹再拜语呢喃，款步芳阶舞绣衫；具道殷勤猫女意，粉资亲制手掺掺。花冠银钏锦为衣，妙舞清歌笑合围；低唱一声金一扣，独留太上古音希。诸峰攒集黛螺青，玉岭如银色独莹；展拓晴云千万里，插天一幅水晶屏。莽葛轻便破碧沙，弄潮如戏语喧哗；依稀佛子乘杯渡，仿佛仙翁泛海槎。倦来高枕树为巢，藤蔓连牵枝格交；栩栩梦回非是蝶，一身幻化类蟏蛸。临溪问渡少，厅，石涧分流远击撞；腰上葫芦头上羽，只身飞过水淙淙。抄阴尺布不堪缝，无褐无衣可耐风；北地乍寒偷射猎，人人尽是鹿皮翁。北辙初回又指南，迢迢原隰未停骖；看山几遍还遵海，阅尽边方烟与岚。成帷成幙逐飞尘，纨裤多缠舆隶身；惯习淫奢无善俗，少年思怕老来贫。二林池运接三林，淡水荥洄咸水深；极目沧波浮海市，一拳真欲笑蹄涔。晓霞散彩覆阴崖，海曲人家逐岸排；烟水几湾帆几幅，顿教风景亿江淮。木冈东峙晓云凝，咫尺阪陀未即登；指点万山分脉处，一峰独秀镇平陵。观音山径几湾环，罗汉门边虎豹关；笑指当年空守戍，但知深谷有乌蛮。陂卧晴沙号七鲲，如环如抱复如蹲；惊涛夜拍殷雷起，远势平吞鹿耳门。龟蛇对峙锁孤城，形势空传统领营；不筑埤头筑海口，为怜安土重纷更。打鼓山头石罅开，悬崖倒拍海潮回；雷声鼎沸浮空翠，万里风樯认影来。矫首南荒欲尽头，影浮拳石小琉球；天回地转渺无际，万水朝宗亘古流。仙山缥渺暗斜曛，石上棋枰旧印纹；沙马矶头人罕到，烂柯樵子话烟云。赤山葱翠漾春烟，沙暖云晴别有天；日午崖边人语沸，村童队队浴汤泉。生成野性气如枭，出没无端血染刀；剥得头颅当户挂，归来轰饮共称豪。


% 原书第 52 页（PDF 第 65 页）
\BookPage{52}{65}
\noindent 内山遥夹外山高，复岭重冈叠翠涛；一带阿猴林下路，须防薮泽有逋逃。八社丁徭力渐纾，番娘饷税蚤捐除；只今宵画辛勤处，谨护官家十万储。生熟番情百种余，半生半熟亦山居；当年户口可知数，尽向鱼鳞册上书。问俗殊方竟未厌，忽教放綷转丹襜；天南水起山穷处，瑯峤云从马首瞻。星轺回处转旌霓，人海无声马不嘶；敢道霜威堪镇俗，长思涛静木城西。

余文仪《续修台湾府志》〈艺文志〉仅录十二首。

又，陈寿祺等《福建通志》（同治十年重刊本）卷百四十〈宦绩〉「国朝「巡台御史夏之芳」传云：夏之芳字荔园，江南高邮州人，雍正癸卯进士。六年，以御史任，兼督学政。不立崖岸，而苞苴自不入其门。按南北二路，民番感悦。主岁科两试，公慎明决，不假督责而士风自化。40

又，《高邮州志》（嘉庆年间重刊本）卷十上，页37-39载：夏之芳，字荔园，号筠庄，赠公绵祚子。雍正癸卯恩科会魁，充内廷教习。乙巳世宗宪皇帝钦命与庶常同应馆试，授翰林院编修大学士。高安、朱文端荐修礼书。出核直隶仓储，策蹇冒风雨，猝至所分县界，往返勾稽，较然顿清。改御史，巡中城，丁未巡视台湾，御赐貂皮、手珠、紫金锭、荷包各数事。旧制台湾学政属台湾道，以芳词臣，优文学，特旨以巡察兼学政事。至则澄叙官方，振兴文教，却私谒，绝苞苴，考课公明，导士子习礼知耻，奏立台字号与乡荐，士气倍奋。按巡南北路，鸡犬不惊，整肃营伍，革惰弁，通事张汉臣[张汉伯]等诱生番杀人，锻炼成狱，以希爵赏，察其枉释诸番，执通事，杖杀界上，画地以杜奸欺，番汉帖然。奉旨留任一年，万里之外，颁赐御书福字，貂皮、哈密瓜以示奖异。庚戌管通州，…再甲寅，巡察山西…乾隆元年丙辰，召掌河南道御史。…后以终养，予告旋里，卒年五十有八。性孝友，应礼闱试时，闻父丧归痛不欲生，奉使台湾，以母老请假省亲，上赐貂皮、手珠、紫金

\vfill
\begin{BookNotes}
\BookNote{40.}{陈寿祺等，《福建通志》卷140，页32b，华文书局影刊同治十年本，台北。（全书五册，未着出版年份。)}
\end{BookNotes}

% 原书第 53 页（PDF 第 66 页）
\BookPage{53}{66}
\noindent 锭、鼻烟壶、火镰荷包，宠赉有加。归田后，别购一宅，迎养萱闱。兄景莱，幼病足疾，曲意承事，终身无闲。弟五人，相继从学，先后入词垣者二人。生平枕图馈史，学综今古，着有《禹贡汇览》、《汉名臣言行录》、《巡台百韵奏疏》一卷。《福建通志》名宦有传。子杏春，辛酉举人，别有传（嘉庆十三年入祀乡贤祠，嘉庆十八年增载）。41

\BookSubheading{（六）刘师恕}

\noindent 《福建通志》卷百之七「国朝职官」页3b，「观风整俗使」条谓:「雍正七年设，八年以后停差。」，刘师恕条云：江南宝应人，康熙庚辰进士，雍正七年任。

《清史稿》卷二百九十一〈列传七十八〉「刘师恕、焦祈年、李徽」传载：刘师恕，字艾堂，江南宝应人…师恕，三十九年进士，选庶吉士，授检讨。累迁国子监祭酒。雍正元年，授贵州布政使。四年，迁通政使，转左副都御史，擢工部侍郎。上以宜兆熊署直隶总督，调师恕礼部，协理总督事。…七年，命师恕以内阁学士充福建观风整俗使。…十一年，师恕以病告，省观风整俗使不复设。乾隆七年，宝应灾，治赈，非贫民例不给。师恕族人诸生洞嗾不得赈者，哄堂罢市。上责师恕不能约束，夺官。南巡迎谒，赐侍读学士衔。二十一年，卒。

\BookSubheading{（七）高元昆}

\noindent 《福建通志》卷百之七「国朝职官」页16b，「盐法道」高元昆条：江南江都人，康熙戊戌年进士，乾隆六年任。

\BookSubheading{（八）杨开鼎}

\noindent 余文仪《续修台湾府志》卷三〈职官志〉「官秩·钦命巡视台湾御史」杨开鼎条（页124）：

\vfill
\begin{BookNotes}
\BookNote{41.}{夏之蓉等，《高邮州志》（据清乾隆48年原修本、嘉庆18年增修、道光25年重校刊本影印），成文书局，1970，台北。}
\end{BookNotes}

% 原书第 54 页（PDF 第 67 页）
\BookPage{54}{67}
杨开鼎：江南甘泉人，己未进士。河南道监察御史。乾隆十四年任。主科试，十五年八月以忧去。

谢金銮等《续修台湾县志》卷二〈政志〉「宪纪·巡察御史」杨开鼎条（页119）记载同上。

又，张子文、郭启传、林伟洲撰《台湾历史人物小传一明清暨日据时期》页632-633载：杨开鼎，字峙塘，又作峙堂、致堂，号玉坡，江南甘泉（今属江苏省）人。乾隆四年（1739）进士。十四年（1749）奉令为巡台御史来台，并兼理学政。在任期间，尝修建台湾府儒学宫，历时一年五个月方建成，并为文以记之。又与士绅议修县学，历时八、九月而成，亦为文以记。又仿其前辈夏之芳取台湾士子之文编为《海天玉尺》、张湄之编《珊枝集》之例，亦取台瀛士子之文，分其门类，评其甲乙等第，汇而编之，名曰《梯瀛集》。十六年八月，因丁忧离台。其后官至湖南衡永郴桂道（张子文）。42

\BookSubheading{(九)陈志泰}

\noindent 王瑛曾《重修凤山县志》卷八〈职官志〉「官秩·本县知县」陈志泰条(页230)：陈志泰：江南甘泉人，举人。乾隆十四年任。

同卷，「宦绩」陈志泰传（页240）云：陈志泰，江南甘泉人，康熙丁酉举人。乾隆十四年由泰宁调知县事，为人廉静耿介，言笑不苟，案牍只字必出己手，绳家人以法，治奸吏蠹役如鹰鹯之逐鸟雀，积弊悉除。讼狱两造，具即听断，民无积滞拘挛之苦。任三年，民无一词、一事枉于情势者。以积劳成疾去，士民思之。43

\vfill
\begin{BookNotes}
\BookNote{42.}{张子文、郭启传、林伟洲撰，《台湾历史人物小传—明清暨日据时期》（国家图书馆出版，民国92年12月）}
\BookNote{43.}{王瑛曾等，《重修凤山县志》（乾隆二十九年，1764)，方豪点校，《台湾丛书·台湾方志汇编》本，国防研究院，1968，台北。}
\end{BookNotes}

% 原书第 55 页（PDF 第 68 页）
\BookPage{55}{68}
\BookSubheading{（十）刘淑}

\noindent 王瑛曾《重修凤山县志》卷八〈职官志〉「官秩·下淡水巡检」刘淑条（页236）：

刘淑：江南甘泉人，供事，乾隆十六年五月任。

\BookSubheading{（十一）丁居信}

\noindent 王瑛曾《重修凤山县志》卷八〈职官志〉「官秩·本县知县」丁居信条（页231):

丁居信：江南仪征人，庶吉士散馆。乾隆二十年九月任。又，同志，卷十二（中）〈艺文志·诗赋·七言律诗〉收丁居信〈凤山任满留别）（页393）：海峤欣逢报政年，愧无异绩纪瑶编。灯残细阅琴堂卷，雨后勤观绿野田。八社莺声迟驿使，半屏日色促行鞭。临岐依恋无他祝，大有频歌入舜弦。

\BookSubheading{（十二）龚宣}

\noindent 余文仪《续修台湾府志》卷十〈武备（二）官秩·澎湖水师协镇〉龚宣条：龚宣：字自耘，江南通州人；壬戌武进士。乾隆二十七年十一月任。

\BookSubheading{（十三）江学举}

\noindent 陈寿祺《福建通志》卷百之七〈国朝职官·盐法道〉（页8b):

江学举：江南如皋人，乾隆三十一年任。

\BookSubheading{（十四）潘正鹏}

\noindent 陈寿祺《福建通志》卷百之七〈国朝职官·库大使〉（页8b）：潘正鹏：江南甘泉人，乾隆三十二年任。


% 原书第 56 页（PDF 第 69 页）
\BookPage{56}{69}
\BookSubheading{（十五）高信可}

\noindent 陈寿祺《福建通志》卷百之七〈国朝职官台湾·嘉义县·佳里兴巡检》（页18a）：

高信可：江南泰州人，乾隆四十一年任。

\BookSubheading{（十六）杨景素}

\noindent 陈寿祺《福建通志》卷百之七〈国朝职官·总督〉(页1b):杨景素：江南甘泉人，乾隆四十三年任。

又粮驿分巡道）（页13b）:杨景素：江南甘泉人，乾隆十七年任。

\BookSubheading{（十七）邵宗尧}

\noindent 陈寿祺《福建通志》卷百之七〈国朝职官·台湾·嘉义县·佳里兴巡检》(页18a）：邵宗尧：江南甘泉人，乾隆五十年任。

\BookSubheading{(十八）丁朝雄}

\noindent 台湾省文献委员会《清史稿台湾资料集辑》列传（之三）第五十九（选自清史稿卷三三五）「丁朝雄」条（页648-649）:

丁朝雄，字伯宜；江苏通州人。自行伍，累擢福建台湾水师副将。乾隆五十一年，以任满，赴部引见。至省城，闻林爽文乱起，朝雄策东港与凤山犄角，爽文所必争；白总督常青，请兵屯东港，断其粮道。常青不能用；遣朝雄还台湾，佐海坛镇总兵郝壮猷讨爽文。五十二年春，壮猷偕朝雄率兵二千余击贼，馘三百、浮二十五。日将暮，贼复来攻；朝雄复杀贼百余。贼始去攻凤山，朝雄乘东门，首诸军入；凤山遂复。黄仕简檄朝雄守安平海口；贼攻府城，朝雄偕知府杨廷桦督兵民力御；贼攻桶盘栈，朝雄为前锋出战，台湾道永福、同知杨廷理率兵民继，复杀贼百余，贼败走。冬，朝雄偕游击倪宾率兵千二百、义


% 原书第 57 页（PDF 第 70 页）
\BookPage{57}{70}
\noindent 民二千余攻东港。东港贼数万，其渠吴豹以海岸浅，度舟不能至，不为备；朝雄遣谍以水注贼𪿫，乘雨至水涨，遣兵民分道登岸杀贼，俘豹。以兵寡，不能克，报常青请益兵；常青令驻港口，护𫗵道。既令攻竹子港，毁贼舟。五十三年春，复攻东港，仍遣谍以水注贼𪿫，督兵攻渡口；贼惊窜，逐三十余里，乃倚山而军。贼夜来犯，朝雄戒勿动；及晓，贼倦，掩击，大破之。爽文遣其徒来援，朝雄筑垒困之；贼溃围出，设伏断其归路，而自将追之，大破贼，遂复东港。福康安上其功，授海坛镇总兵。既，福康安劾柴大纪受陋规，言朝雄为安平协副将时亦有此；当夺职，戍军台。上以朝雄攻东港，战有功；命留任。林剃舵、林明灼者，海盗渠也·五十四年，朝雄巡洋至汜澳，破盗巢，得ソ舵等。而明灼拒，杀参将张殿魁，上责总督伍拉纳；伍拉纳以属朝雄·督舟师出海，遇诸大麦洋，俟其近，发大𪿫毙数酋；明灼穷蹙，跃入海，官军钩致。俘以归·五十五年，追论朝雄在台湾失察天地会邪教，当夺职；上咨伍拉纳朝雄在官状·伍拉纳言朝雄督水师捕盗有劳；命还任·五十八年，摄水师提督。五十九年，入觐；至清江浦，病笃，乞罢归，卒于上海舟中。44

\BookSubheading{（十九）吴文镕}

\noindent 陈寿祺《福建通志》卷百之七〈国朝职官·巡抚〉吴文镕条（页3a):吴文镕，江苏仪征人，嘉庆己卯进士，十九年以刑部侍郎任。

又，英杰等《续纂扬州府志》（同治十三年（1874年）刊本）卷九〈人物志

\BookSubheading{（一）)（页35a-37a）:}

\noindent 吴文镕，字甄甫，仪征人。嘉庆十八年举人，二十四年进士。改庶吉士，授编修。道光…十九年，授福建巡抚两署闽浙总督。

\BookSubheading{（二十）吴篪}

\noindent 陈寿祺《福建通志》卷百之七〈国朝职官·粮驿道〉吴篪条（页14a):吴箎，江苏如皋监生，道光元年任。

\vfill
\begin{BookNotes}
\BookNote{44.}{台湾省文献委员会，《清史稿台湾资料集辑》，1994，南投。}
\end{BookNotes}

% 原书第 58 页（PDF 第 71 页）
\BookPage{58}{71}
\BookSubheading{（二十一）易金杓}

\noindent 陈寿祺《福建通志》卷百十七〈国朝职官·台湾〉：

猫雾辣巡检「江苏仪征监生，道光十年任」。（页21b)

新庄县丞「江苏仪征监生，道光十七年任」。(页5a)

头围县丞「江苏仪征监生，道光十九年任」。（页6a)

《台湾文献丛刊》第一百六十种《噶玛兰厅志》卷二(中）〈职官·官秩·头围县丞〉：易金杓：江苏仪征人，由监生南安县丞，十九年九月十四日任。

\BookSubheading{（二十二）高鸿飞}

\noindent 英杰《续纂扬州府志》卷七〈人物·忠义〉，高鸿飞传（页2a-2b）45:高鸿飞，字南卿，号伯鸾，高邮人。幼颖异，勤苦力学。为文，援笔立成，尤工词赋。教授生徒，多所成就。道光五年，拔贡生。十二年举人。二十一年进士。由翰林院庶吉士改知县，选福建福鼎县。首禁火葬陋习。二十五年调署晋江县。遇械鬬案，惩办尤严。嗣补台湾彰化县继委署凤山县彰化葫芦墩。有漳泉人买卖殴鬬，匪徒乘机思逞，村民多迁徙。鸿飞会营，驰赴查拿，诛首祸者，余皆敛戢，地方安堵。先是以协拿洋盗出力，加五品衔；至是以同知补用。咸丰二年，调署台湾县。四月二十八日，突闻凤山县属匪徒聚众台邑蠢应，因会营亲往剿办。出城至湾里街，侦报贼势已炽。外无援兵，贼直前攻扑，兵勇不敌，遂以力竭阵亡。

\vfill
\begin{BookNotes}
\BookNote{45.}{根据中国新闻网2009年3月11日的报导，新发现高鸿飞长生牌。高约60厘米、宽约11厘米，木质牌位竖排阳刻三行镏金字。正中共有32个字—敕授儒林郎知福鼎县事前翰林院庶吉士高公名鸿飞字伯鸾号南卿长生禄位。左边共22个字—辛丑恩科进士完成宾兴考棚诸盛举邑绅士感德不忘。右边共24个字一公系江苏扬州府高邮州人由拔贡生中式道光壬辰科举人。以上资料与《续纂扬州府志》记载有出入。附上长生牌位图片。资料来源：扬州网转贴扬州晚报2009年3月10日报导网址:http://www.yznews.com.cn/news/2009-03/10/content\_2587392.htm}
\end{BookNotes}

% 原书第 59 页（PDF 第 72 页）
\BookPage{59}{72}
\noindent 恤予世职，入祀台湾名宦祠、京师昭忠祠，并于死事地方建立专祠。子蓉镜，同治九年举人。

\BookSubheading{（二十三）徐宗干}

\noindent 台湾省文献委员会，《台湾省通志》〈人物志·宦绩篇·清代〉徐宗干传(页110a-110b）：徐宗干，字树人，江苏南通州人。嘉庆二十五年进士。道光元年，任山东曲阜知县；除武城，调泰安，在任十年，有政声，既迁高唐知州。道光十七年，潍县教匪马刚作乱，从巡抚经额布剿擒之，即地诛之，众心以定。迁济宁直隶州，全乡民因濬河互斗，致殴官，伤胥吏，势汹汹。宗干往谕，使解散；殴官民自首，大吏欲置重典，宗干以为「民畏水患，非与官敌；聚众本沮工，殴官非本意」；力争，仅戌为首者七人。旋署衮州知府，修滋阳河堤。道光二十二年，擢四川保宁知府，兼署川北道。二十四年，授汀漳龙道，民多途控；宗干悉摘传原、被质询，是非立判。乡民素悍辄出掠财械斗，宗干召其家长者至，予酒食劝解，许自新；并籍其中壮丁姓名为乡勇，隶之官，由是暴民敛迹。漳俗迎神报赛，藉是箕敛赀财，举国若狂；宗干悉示禁之，期存古逐疫报功之意而止。道光二十五年，丁母忧，去官。服阕，二十八年四月，起授分巡台湾道，时姚莹方去职，凡所规画，多继成之。宗干为治，循名核实；振兴文教，汲汲以育才为务。台湾遭英人窥扰之后，士民蓄愤，自立乡约，禁不与贸易。宗干亦着防夷论，倡兵民防堵，清内奸。壮士卒，以联庄缉匪为名，实则以防英人登岸也。当是时，绿营废弛，班兵多宿民家，挟械以嬉。宗干移镇管束，改建营房处之。兵民始分。又议改澎湖募兵，变通船政，清理人犯，语多可行。水沙连六社番久请内附，而廷议以险远为难，照旧封禁。宗干上书总督，请援乾隆五十三年之例，先设屯丁，以便管理。从之；其后遂设官焉。咸丰三年，凤山林恭起事，陷县治，攻府城，复援噶玛兰厅，宗干与绅民守御，命知县郑元杰以兵平之。四年，升福建按察使；为巡抚王懿德所劾，解任，旋召赴京，命往河南襄办军务。六年，复命赴安徽。七年，授浙江按察使，迁布政使。十年，江苏团练大臣庞钟璐请以宗干办理通、泰诸州县团练。同治元年四月，擢福建巡抚，彰化戴潮春事起，全台俱扰；而福建上游，军务复急。省议


% 原书第 60 页（PDF 第 73 页）
\BookPage{60}{73}
\begin{BookQuote}
颇不以台湾为意，宗干独顾念焉。即檄前署台湾镇曾玉明渡台，又奏简丁曰健为分巡台湾道，会办军务，次第平之。五年，病卒。闽浙总督左宗棠偕将军英桂奏：「宗干循良着闻，居官廉惠得民，所至有声。」优诏褒恤，谥清惠。同治七年，入祀福建名宦祠。
\end{BookQuote}

《清史稿》、《漳州府志》、《光绪台湾通志》、《台湾通史》、《台湾省通志》等书，皆为徐宗干设传，不赘录。台湾银行经济研究室《台湾文献丛刊》搜集徐着《斯未信斋文集》，取其与台湾相关部分为《斯未信斋文编》列为第八十七种，另有《斯未信斋杂录》列为第九十三种。徐宗干所辑《治台必告录》则列为第十七种。

\BookSubheading{（二十四）厉恩官}

英杰《续纂扬州府志》卷九〈人物志〉，厉恩官传（页47b-48a）：

\begin{BookQuote}
厉恩官，字锡功，号砚秋。仪征人。道光十二年顺天举人。二十年进士。改庶常授编修。二十三年分校顺天乡闱。二十四年典试河南，补山东道御史，掌京畿道。京察外用任山东衮沂曹济道，升按察使。咸丰四年，粤匪窜踞连镇，奉檄驰赴陵县防堵。在事年余，守御严密，边境获安。以功赏戴花翎，升布政使。转太常寺卿，迁宗人府府丞，署兵刑两部侍郎。十一年视学福建兼署巡抚，差竣回京，遘疾，卒，年五十五。弟恩绶，字慎功。诸生援例分发山西官荥水县知县，卒于任。
\end{BookQuote}

\BookSubheading{（二十五）王凯泰}

台湾省文献委员会，《台湾省通志》卷七〈人物志·宦绩篇·清代〉王凯泰传（页134a）：

\begin{BookQuote}
王凯泰，字补帆，江苏宝应人。道光三十年进士。同治九年，以广东布政使迁福建巡抚，甫下车，为六言韵语；劝士以立品、励学；劝民以息讼、止斗、禁火葬、戒溺女。同治十三年，观京师。未至，适台湾有事，诏回任。时沈葆桢方督台湾海防，奏请移巡抚其地。又奏抚臣王凯泰励行清苦，用心缜密；开山
\end{BookQuote}

抚番诸大役，宜以任之。凯泰亦奏，事难遥断，请先往履勘奏牍。光绪元年五


% 原书第 61 页（PDF 第 74 页）
\BookPage{61}{74}
\noindent 月渡台；台民嗜博；且无长幼皆食外国烟；又锢婢，至老不使嫁；凯泰作五言百句歌以劝化之。居台五月，力筹整顿，并兼顾省大局，次第以事入告；而积劳受瘴疠，十月内渡，卒于官，年五十有三。赠太子少保衔，视总督例。赐恤，谥文勤。命于福建省城、台湾府城建祠焉。

王凯泰协助沈葆桢在台推行新政，两相契合，《沈文肃公牍》中，收录文肃致文文勤书凡三十函，不赘引46。

王凯泰去世后，除了史馆立传之外，闽浙总督何璟等也奏请将王凯泰入祀名宦，其文曰：闽督何（璟）奏已故大员恳入祀名宦折（正月二十七日京报)闽浙总督兼署福州将军臣何璟、署理福建巡抚调补山西布政使臣葆亨跪奏：为已故大员功德在民，吁恳入祀「名宦」以光盛典；恭折仰祈圣鉴事。窃维圣世旌贤，百代树人臣之表；明禋报德，千秋隆秩祀之仪；所以彰暗循良、懋昭忠盖也。兹据绅士前湖北巡抚郭柏荫等呈称:「原任福建巡抚、追赠太子少保王凯泰，起家词馆，从事戎轩；陈皋浙江，开藩粤口；历官政绩，遐迩周知。泊其秉节入闽，典武阐而鉴空衡平，司秋试则风清弊绝。通经致用，拓广厦以庇寒儒；整俗观风，怀好音而集泮水。躬行俭德，首毁淫词。既而奉命渡台，备尝艰苦；随行仆众，司行李者两丁；手治官书，对青灯于午夜。恶牛鬼蛇神之说，无取狂夫；杜城狐社鼠之奸，先绳悍卒。为百姓力除鸩毒，烟禁特申；悬三尺以戒猪奴，博徒知儆。身婴瘴疠，始返重洋；病入膏育，竟成朝露！囊中长物，惟有图书；身后去思，徧于妇孺，特颁世赏，子孙并锡簪缨；予建专祠，血食常沾俎豆；易名彰其文德，赐血以表勤劬；固以恩叠颁，重泉知感矣。柏荫等久居编蔀，同志甘棠。岘山望羊傅之碑，依然堕泪；颍水溯雍奴之治，无自扳辕。墓木犹新，人其宛在；涧苹用荐，礼亦宜之。开具事迹，佥恳转详奏请入祀名宦祠，以顺舆情。至故官子弟，并无现官九卿及曾任九卿。」等情。由署布政使李明墀会同署按察使叶永元照造事迹清册，详情具奏前来。臣等伏查原任福建巡抚、追赠太子少保王凯泰，自翰林起家，历官司道，游膺

\vfill
\begin{BookNotes}
\BookNote{46.}{沈葆桢，《沈文肃公牍》，台湾省文献委员会，1998，台中。}
\end{BookNotes}

% 原书第 62 页（PDF 第 75 页）
\BookPage{62}{75}
\begin{BookQuote}
疆寄；历中外几二十年。所至惟以正己率属为心，而于地方利弊尤所留意；莫不乘时因革，抉剔靡遗。至今奉为楷模，遵循不易。台湾之役，开山抚番、征兵筹饷，殚竭心力，瘴疠交攻；染疾卧床，犹伏枕办事，强作支持。殆疴渐沈，始从僚属之劝，扶病内渡；而膏肓莫救，且不起矣。身后囊橐萧然，无殊寒畯。柩回原籍之日，阖城士庶为之泣下；至今父老谈其遗事，莫不感慨歔欷，声泪俱堕。盖其遗泽入人之深，历久弗渝有如此者。前蒙圣恩赐恤、予谥，建立专祠，荫及后裔；是其官声政绩，久邀圣明洞鉴之中。现据绅士前湖北巡抚郭柏荫等呈请入祀「名宦」以伸遗爱，良由爱戴之深，不能自已。第督、抚入祀「名宦」，例应具题；今该绅等呈请具奏，核与定例不符。惟查原任福建巡抚徐宗干入祀福建名宦祠，奏奉谕旨允准有案，相应援照改题为奏。合无仰恳天恩，俯准原任福建巡抚、追赠太子少保王凯泰入祀福建名宦祠，以顺舆情而褒忠盖。47
\end{BookQuote}

以上二十五人中，担任福建总督、巡抚、盐法道、库大使等统辖官者九人；巡台御史两人；台湾道以下知县、县丞、巡检等文职者九人、武职五人，就一府的人才而言，扬州与台湾的关系不可谓不密切。

五结论

康熙二十三年（1684），设台湾府，下辖台湾、诸罗、凤山三县。在传统地方志中，「星野说」占绝对优势，台湾地志，皆有星野一门，在星野说中，福建属于扬州；台湾既为福建一府，则台湾亦属于扬州。台湾方志，皆冠诸篇首，咸无异辞。

清初，台湾盛产米、糖、麻、豆、鹿肉、鹿皮，行销各国及大陆沿海各省，尤其长江三角洲各城市；回程时则载运各地特产销台。上海、苏州、镇江皆入记载，唯扬州虽与三城近在密迩却未见诸记载，或可就扬州地方文献或各商家私藏档案中求之。

\vfill
\begin{BookNotes}
\BookNote{47.}{〈闽督何奏已故大员恳入祀名宦折〉，《申报》光绪四年二月二十日（1878）条，页768769。}
\end{BookNotes}

% 原书第 63 页（PDF 第 76 页）
\BookPage{63}{76}
根据日治时期（1927）「台湾在籍汉民族乡贯别调查」的报告，1925年底设籍在台湾的汉人有3751600人，80\%来自地近台湾的闽南（泉州府44.8\%、漳州府35.2\%）；1.6\%来自闽西（汀州府和龙岩州）；1.5\%来自闽北（福州府、兴化府和永春州）；15.6\%来自粤东（潮州府、惠州府和嘉应州）。「其他」项下只有48900人，占1.3\%。48

笔者无法得知当年日本人进行调查时，「其他」项下包含哪些省份和府、州、县。若依前述闽、粤两省之部分府、州、县即占98.7\%，假设两省之其余府、州、又占1.3\%之98.7\%，则其他各省人民尚不及64人。若从宽估算，其他各省人民占「其他」项下之30\%，亦不过1467人。

就前述台湾汉人之祖籍构成观察，台湾可说是相当封闭的社会，极少闽粤两省以外之各省人民至此定居。若然，纵使有扬州人在台湾定居，也只是少数；另一方面，扬州是盐业极盛的消费型城市，吸纳大量外地盐商与劳动人口49，亦可能有极少数台湾人迁徙扬州。以上推论，要想得到确切的证据，应属可遇而不可求。

乾隆五十一年（1786）台湾爆发林爽文事变，全台大乱，惊动全国，各省商人捐银挹住政府财政，两淮商人捐银两百万两，犹胜于其他各地商人捐银之总和，通州人丁朝雄于斯役战绩卓着，授海坛镇总兵，有功于台湾。

甘泉人江藩（着有《国朝汉学师承记》、《宋学渊源记》）因大学士王杰之荐，受乾隆帝赏识「谕召对圆明园」，「值林爽文陷台湾报至，遂辍，人惜其遇」50，可视为林爽文事件之扬州受害人。

台湾乃新垦之地，移垦社会民风剽悍、桀骜难治。高邮人高鸿飞，历任台湾、彰化、凤山三县知县，咸丰二年（1852）因剿匪而阵亡于湾里街（今台南市南区)。

\vfill
\begin{BookNotes}
\BookNote{48.}{台湾总督府官房调查课，《台湾在籍汉民族乡贯别调查》页4-5，昭和三年(1928)，台北。}
\BookNote{49.}{参见冯尔康编，《扬州研究—江都陈轶群先生百龄冥诞纪录论文集》，陈捷先出版，民国八十五年(1996)，台北。所收：徐明德，〈论唐代扬州国际大港的繁荣与历史地位〉，页139-177。王振忠，〈明清徽商与扬州城市文化的特征与地位〉，页489-520。}
\BookNote{50.}{英杰等，《续纂扬州府志》卷13〈人物志·文苑·江藩传〉（页9b-10b)。}
\end{BookNotes}

% 原书第 64 页（PDF 第 77 页）
\BookPage{64}{77}
通州人徐宗干，道光二十八（1848）任台湾道，励精图治，所辑《治台必告录》，汇整百余年治台、知台之智慧而成，不仅清末官吏奉为圭臬，抑为今日研究台湾史之宝典。同治元年（1862）任福建巡抚，戡平仅次于林爽文事件之戴潮春事件，五年，卒于官。

宝应人王凯泰，同治九年任闽抚，十三年（1874）日人侵台，沈葆桢督办台湾海防，王凯泰于光绪元年（1875）五月渡台，负责善后事宜，并亲理开山抚番，积劳成疾，十月内渡，卒于官。

以上扬州人，皆为台湾鞠躬尽瘁，死而后已，令人感佩。

四十余年前，受业于江都夫子陈老师捷先先生，受恩独厚。读史，偶及于扬州辄录存之。扬州所属州县文献及世家大族家藏档案均未寓目，本文谨将多年所辑录史料汇聚成篇，抛砖引玉也，聊表谢意于万一而已。
